% GENERATED FILE. Edit canonical lessons, then run npm run generate:books.
% canonical-source-hash: fnv1a64:d1e40abd18092cc2
% canonical-lessons: ZH-C386-ershou, ZH-C386-biancheng, ZH-W386-bian-change, ZH-C386-xianzhi, ZH-C386-huasuan, ZH-C386-biantian

\chapter{Second-Hand}
\label{ch:zh-second-hand}
\hlchaptermodality{386}

\begin{chapteropening}
\textbf{By the end of this chapter:} \emph{I can talk about second-hand things, how something turned into something else, unused things, good value and the weather turning.}

Complete the last lesson of chapter 386: I can talk about second-hand things, how something turned into something else, unused things, good value and the weather turning.
\end{chapteropening}

\section[\texorpdfstring{\zh{二手} (èrshǒu)}{二手 (èrshǒu)}]{\zh{二手} (èrshǒu) — second-hand}

\label{lesson:ZH-C386-ershou}



\begin{quote}
Before the new one: say the Mandarin for to attach importance to, then the Mandarin for a way of doing things.
\end{quote}

\subsection*{You'll want to know: \zh{二手}}
\textbf{\zh{二手}} — \emph{èrshǒu} — \textquotedblleft{}second-hand\textquotedblright{}.

\begin{quote}
\emph{wǒ mǎi le yí ge èrshǒu shǒujī yīnwèi hěn piányi} — I bought a second-hand phone because it was cheap.
\end{quote}

\begin{cousinweb}
\textbf{\zh{二手}} (\emph{èrshǒu}) is built from \zh{二} 'two' and \zh{手} 'a hand', both inherited from Middle Chinese, \emph{nyijH} and \emph{syuwX} (Guangyun, in Baxter's transcription): through a second pair of hands.
\end{cousinweb}

\begin{grammarlens}[title={What you've built}]
One more word for giving reasons.
\end{grammarlens}

\subsection*{Your turn}
\begin{itemize}
\raggedright
  \item \emph{Say it:} \emph{èrshǒu}
  \item \emph{Say it:} \emph{èrshǒu}, once more
  \item \emph{Say it:} \emph{wǒ mǎi le yí ge èrshǒu shǒujī yīnwèi hěn piányi} — I bought a second-hand phone because it was cheap.
  \item \emph{Recall:} say the Mandarin for to attach importance to, then the Mandarin for a way of doing things, then say \emph{èrshǒu} again
\end{itemize}

\begin{tcolorbox}[breakable,colback=teal!4,colframe=teal!35!black,title={Before you move on}]
What does \zh{二手} mean? (\textquotedblleft{}Second-hand\textquotedblright{}.) Say it once more.
\end{tcolorbox}
\section[\texorpdfstring{\zh{变成} (biànchéng)}{变成 (biànchéng)}]{\zh{变成} (biànchéng) — to turn into}

\label{lesson:ZH-C386-biancheng}



\begin{quote}
Before the new one: say the Mandarin for a way of doing things, then the Mandarin for second-hand.
\end{quote}

\subsection*{You'll want to know: \zh{变成}}
\textbf{\zh{变成}} — \emph{biànchéng} — \textquotedblleft{}to turn into\textquotedblright{}.

\begin{quote}
\textbf{\zh{那家饭店变成了书店}} \emph{nà jiā fàndiàn biànchéng le shūdiàn} — That restaurant has turned into a bookshop.
\end{quote}

The new character, in pieces: \zh{变} — eight strokes in two pieces: on top a dot and a bar over four short strokes; at the foot \zh{又}, two strokes. The next lesson puts it in your hand.

\begin{cousinweb}
\textbf{\zh{变成}} (\emph{biànchéng}) is built from \zh{变} 'to change' and \zh{成} 'to complete', both inherited from Middle Chinese, \emph{pjenH} and \emph{dzyeng} (Guangyun, in Baxter's transcription), which glosses the first 'to transform'.
\end{cousinweb}

\begin{grammarlens}[title={What you've built}]
One more word for giving reasons.
\end{grammarlens}

\subsection*{Your turn}
\begin{itemize}
\raggedright
  \item \emph{Say it:} \emph{biànchéng}
  \item \emph{Say it:} \emph{biànchéng}, once more
  \item \emph{Say it:} \emph{nà jiā fàndiàn biànchéng le shūdiàn} — That restaurant has turned into a bookshop.
  \item \emph{Recall:} say the Mandarin for a way of doing things, then the Mandarin for second-hand, then say \emph{biànchéng} again
\end{itemize}

\begin{tcolorbox}[breakable,colback=teal!4,colframe=teal!35!black,title={Before you move on}]
What does \zh{变成} mean? (\textquotedblleft{}To turn into\textquotedblright{}.) Say it once more.
\end{tcolorbox}
\section[\texorpdfstring{\zh{变} (biàn)}{变 (biàn)}]{\zh{变} — a dot and a bar over four short strokes, then \zh{又}}

\label{lesson:ZH-W386-bian-change}



\begin{quote}
The word first, then the character inside it.

\begin{itemize}
\raggedright
  \item \emph{Recall:} say \emph{biànchéng} — \textbf{R1}, one lesson back, before it decays
\end{itemize}

\begin{quote}
\zh{文} \zh{又}
\end{quote}

\emph{Read it:} each one aloud, with nothing said first — \emph{wén}, \emph{yòu}
\end{quote}

\subsection*{Script — a dot and a bar over four short strokes, then \zh{又}}
\hlblockfigure{\detokenize{figures/ZH-W386-bian-change-filmstrip.pdf}}{How \zh{变} is written, stroke by stroke}

\begin{quote}
\zh{变}
\end{quote}

\textbf{Two pieces, top before bottom.}

\textbf{The top piece, six strokes.} A dot, then a long bar under it, as at the top of \zh{文}. Two short verticals under the bar, the left one first. Then a short stroke falling to the left beside the left vertical, and a dot falling to the right beside the right one.

\textbf{The bottom piece, \zh{又}, two strokes,} as you wrote it: a bar that turns and falls to the lower left, then a long stroke falling to the right across it.

6 + 2 = eight.

\subsection*{Your turn}
\begin{itemize}
\raggedright
  \item \emph{Write it:} \zh{变}, all eight, with the model in front of you
  \item \emph{Cover:} write it again from memory, and count the pieces, not the strokes
\end{itemize}

\begin{tcolorbox}[breakable,colback=teal!4,colframe=teal!35!black,title={Before you move on}]
How many strokes in \zh{变}? (\textbf{Eight}.) How does it start? (\textbf{A dot and a bar, as in \zh{文}}.) What is the bottom piece? (\textbf{\zh{又}, two strokes}.)

Source: \href{https://raw.githubusercontent.com/chanind/hanzi-writer-data/68d10a4b21150cae5e1ebbd223eed289cf32d90c/data/\%E5\%8F\%98.json}{Hanzi Writer Data \zh{变}, pinned stroke paths and medians}.
\end{tcolorbox}
\section[xiánzhì]{xiánzhì — unused; lying idle}

\label{lesson:ZH-C386-xianzhi}



\begin{quote}
Before the new one: say the Mandarin for second-hand, then the Mandarin for to turn into.
\end{quote}

\subsection*{You'll want to know: xiánzhì}
\emph{xiánzhì} — \textquotedblleft{}unused; lying idle\textquotedblright{}.

\begin{quote}
\emph{jiā lǐ xiánzhì de dōngxi kěyǐ sòng gěi péngyou} — Things you don't use at home can go to friends.
\end{quote}

Its characters come later; for now it is a word for the ear.

\begin{cousinweb}
\emph{xiánzhì} is built from \emph{xián} 'idle' and \emph{zhì} 'to put', both inherited from Middle Chinese, \emph{hean} and \emph{triH} (Guangyun, in Baxter's transcription), which glosses the first 'leisure' and the second 'to set down'.
\end{cousinweb}

\begin{grammarlens}[title={What you've built}]
One more word for giving reasons.
\end{grammarlens}

\subsection*{Your turn}
\begin{itemize}
\raggedright
  \item \emph{Say it:} \emph{xiánzhì}
  \item \emph{Say it:} \emph{xiánzhì}, once more
  \item \emph{Say it:} \emph{jiā lǐ xiánzhì de dōngxi kěyǐ sòng gěi péngyou} — Things you don't use at home can go to friends.
  \item \emph{Recall:} say the Mandarin for second-hand, then the Mandarin for to turn into, then say \emph{xiánzhì} again
\end{itemize}

\begin{tcolorbox}[breakable,colback=teal!4,colframe=teal!35!black,title={Before you move on}]
What does xiánzhì mean? (\textquotedblleft{}Unused; lying idle\textquotedblright{}.) Say it once more.
\end{tcolorbox}
\section[huásuàn]{huásuàn — good value; worth it}

\label{lesson:ZH-C386-huasuan}



\begin{quote}
Before the new one: say the Mandarin for to turn into, then the Mandarin for sitting unused at home.
\end{quote}

\subsection*{You'll want to know: huásuàn}
\emph{huásuàn} — \textquotedblleft{}good value; worth it\textquotedblright{}.

\begin{quote}
\emph{wǒ juéde èrshǒu chē bǐ xīn chē huásuàn} — I think a second-hand car is better value than a new one.
\end{quote}

Its characters come later; for now it is a word for the ear.

\begin{cousinweb}
\emph{huásuàn} is built from \emph{huá} and \emph{suàn} 'to count'. The second is inherited from Middle Chinese \emph{swanX} (Guangyun, in Baxter's transcription), glossed 'the number of things'; the dictionary glosses the first 'to row a boat', and how it came into this word is not settled.
\end{cousinweb}

\begin{grammarlens}[title={What you've built}]
One more word for giving reasons.
\end{grammarlens}

\subsection*{Your turn}
\begin{itemize}
\raggedright
  \item \emph{Say it:} \emph{huásuàn}
  \item \emph{Say it:} \emph{huásuàn}, once more
  \item \emph{Say it:} \emph{wǒ juéde èrshǒu chē bǐ xīn chē huásuàn} — I think a second-hand car is better value than a new one.
  \item \emph{Recall:} say the Mandarin for to turn into, then the Mandarin for sitting unused at home, then say \emph{huásuàn} again
\end{itemize}

\begin{tcolorbox}[breakable,colback=teal!4,colframe=teal!35!black,title={Before you move on}]
What does huásuàn mean? (\textquotedblleft{}Good value; worth it\textquotedblright{}.) Say it once more.
\end{tcolorbox}
\section[\texorpdfstring{\zh{变天} (biàntiān)}{变天 (biàntiān)}]{\zh{变天} (biàntiān) — to turn bad (of the weather)}

\label{lesson:ZH-C386-biantian}



\begin{quote}
Before the new one: say the Mandarin for sitting unused at home, then the Mandarin for good value for money.
\end{quote}

\subsection*{You'll want to know: \zh{变天}}
\textbf{\zh{变天}} — \emph{biàntiān} — \textquotedblleft{}to turn bad (of the weather)\textquotedblright{}.

\begin{quote}
\textbf{\zh{要变天了你最好早回家}} \emph{yào biàntiān le nǐ zuìhǎo zǎo huí jiā} — The weather's about to turn; you'd better get home early.
\end{quote}

\begin{cousinweb}
\textbf{\zh{变天}} (\emph{biàntiān}) is built from \zh{变} 'to change' and \zh{天} 'the sky', both inherited from Middle Chinese, \emph{pjenH} and \emph{then} (Guangyun, in Baxter's transcription).
\end{cousinweb}

\begin{grammarlens}[title={What you've built}]
One more word for giving reasons.
\end{grammarlens}

\subsection*{Your turn}
\begin{itemize}
\raggedright
  \item \emph{Say it:} \emph{biàntiān}
  \item \emph{Say it:} \emph{biàntiān}, once more
  \item \emph{Say it:} \emph{yào biàntiān le nǐ zuìhǎo zǎo huí jiā} — The weather's about to turn; you'd better get home early.
  \item \emph{Recall:} say the Mandarin for sitting unused at home, then the Mandarin for good value for money, then say \emph{biàntiān} again
\end{itemize}

\begin{tcolorbox}[breakable,colback=teal!4,colframe=teal!35!black,title={Before you move on}]
What does \zh{变天} mean? (\textquotedblleft{}To turn bad (of the weather)\textquotedblright{}.) Say it once more.
\end{tcolorbox}
